\documentclass[aspectratio=169]{beamer}
\usetheme{Madrid}
% \usecolortheme{beaver}

\usepackage{tikz}
\definecolor{beaver}{RGB}{146,111,91}
\usetikzlibrary{arrows.meta,positioning}

\setbeamertemplate{itemize items}[triangle]
\setbeamertemplate{enumerate items}[circle]
\setbeamersize{text margin left=1cm,text margin right=1cm}

\title{Towards a Quantum Compiler}
\subtitle{Joint PhD Project Elevator Pitch}
\author{Henri \& Clemens}
\date{\today}

\begin{document}

\begin{frame}[plain]
  \titlepage
\end{frame}

\begin{frame}[t]{Vision}
  \begin{columns}[T,onlytextwidth]
    \column{0.58\textwidth}
      \begin{block}{North Star}
        \begin{itemize}
          \item Build a compiler that closes the gap between high-level quantum descriptions and efficient circuits.
          \item Use ZXW calculus as the reasoning layer that keeps simplifications provably sound.
        \end{itemize}
      \end{block}
      \vspace{0.25cm}
      \begin{itemize}
        \item Contribute tooling the community can adopt quickly without rewriting their stack.
      \end{itemize}
    \column{0.38\textwidth}
      \begin{alertblock}{What Success Looks Like}
        \begin{itemize}
          \item QUIZX gains W-spider automation that researchers rely on daily.
          \item Our compiler prototypes motivate new hardware experiments within months.
        \end{itemize}
      \end{alertblock}
  \end{columns}
\end{frame}

\begin{frame}[t]{Compiler Flow at a Glance}
  \centering
  \tikzset{
    stage/.style={rounded corners, draw=beaver!70!black, fill=beaver!20, minimum width=3.1cm, minimum height=1.1cm, align=center},
    note/.style={font=\small, align=center}
  }
  \resizebox{0.96\linewidth}{!}{%
    \begin{tikzpicture}[node distance=2.2cm, every node/.style={font=\small}, >=Latex, thick]
      \node[stage] (model) {High-level\\quantum model};
      \node[stage, right=of model] (zxw) {ZXW diagrams\\\& simplification};
      \node[stage, right=of zxw] (results) {Optimized circuits\\+ numeric summaries};
      \node[stage, right=of results] (hardware) {Target hardware\\constraints};
      \node[stage, below=1.1cm of zxw, minimum width=3.6cm] (feedback) {Classical feedback:\\flow, gflow, heuristics};

      \draw[->] (model) -- (zxw);
      \draw[->] (zxw) -- (results);
      \draw[->] (results) -- (hardware);
      \draw[->] (hardware.south) |- ++(0,-0.5) -| (feedback.east);
      \draw[->] (feedback.north) -- (zxw.south);

      \node[note, above=0.25cm of zxw] {ZXW calculus keeps\\rewrites certifiable};
    \end{tikzpicture}%
  }
  \vspace{0.35cm}
  \begin{columns}[T,onlytextwidth]
    \column{0.48\textwidth}
      \begin{block}{What Enters the Loop}
        \small High-level models, Pauli-sum Hamiltonians, and hardware constraints surface together so every rewrite stays grounded in the target experiment.
      \end{block}
    \column{0.48\textwidth}
      \begin{exampleblock}{What We Get Out}
        \small Numerical benchmarks plus hardware-ready circuits, with classical feedback guiding the next simplification pass.
      \end{exampleblock}
  \end{columns}
\end{frame}

\begin{frame}[t]{Impact in Three Moves}
  \begin{enumerate}[<+->]
    \item Compile classically tractable circuits to numerical predictions for threshold analysis.
    \item Shrink circuit depth and qubit counts before experiments hit scarce hardware.
    \item Extract circuits that match restricted gate sets, reducing calibration overhead.
  \end{enumerate}
  \vspace{0.3cm}
  \begin{exampleblock}{Bottom Line}
    Faster iteration loops for both theorists and experimentalists exploring quantum processes.
  \end{exampleblock}
\end{frame}

\begin{frame}[t]{Technical Focus}
  \begin{columns}[T,onlytextwidth]
    \column{0.48\textwidth}
      \begin{block}{QUIZX Enhancements}
        \begin{itemize}
          \item Invent and implement W-spider rewrite routines.
          \item Add scalable notation so large diagrams remain legible.
        \end{itemize}
      \end{block}
    \column{0.48\textwidth}
      \begin{exampleblock}{Research Threads}
        \begin{itemize}
          \item Classical feedback for QEC diagrams with mixed spider types.
          \item Flow and gflow driven circuit extraction heuristics.
        \end{itemize}
      \end{exampleblock}
  \end{columns}
\end{frame}

\begin{frame}[t]{Immediate Experiments}
  \begin{columns}[T,onlytextwidth]
    \column{0.32\textwidth}
      \begin{block}{Week 0--2}
        \small Create ZX diagrams from Pauli-sum Hamiltonians and benchmark baseline simplifications.
      \end{block}
    \column{0.32\textwidth}
      \begin{block}{Week 3--5}
        \small Explore rewrite rules involving Z boxes, \(\pi\) rotations, and W spiders; catalogue promising patterns.
      \end{block}
    \column{0.32\textwidth}
      \begin{block}{Week 6+}
        \small Upstream minimal yet compelling QUIZX contributions and automate regression checks.
      \end{block}
  \end{columns}
  \vspace{0.3cm}
  \begin{alertblock}{Support We Need}
    Quick feedback on simplification heuristics and access to representative QEC diagrams.
  \end{alertblock}
\end{frame}

\begin{frame}[plain]
  \centering
  \vfill
  {\Large Join us to stress-test W-spider rules, co-design benchmarks, and land the first QUIZX pull request together.}
  \vfill
\end{frame}

\end{document}
